\reviewercomment{2. Evidence for community-level selection

The manuscript interprets dominance of one parental community and correlated persistence of its taxa as evidence for community-level selection. We find this interpretation interesting and potentially insightful, but it is not uniquely diagnostic on its own.

Similar patterns could arise from:
- shared environmental filtering
- correlated traits within parental communities
- environmental modification (e.g. pH tolerance)

In particular, the nutrient-rich regime appears consistent with strong environmental filtering imposed by a small number of dominant taxa, which may not require invoking selection acting at the level of the community as a unit.

We therefore recommend clarifying the interpretation of "community-level selection," and explicitly discussing alternative mechanisms that could generate similar patterns. This distinction has been discussed in previous work on community coalescence and microbial assembly [Mansour et al., 2018, Rillig et al., 2015].

Relatedly, we suggest clarifying the ecological meaning of "failed invasion." At present it is primarily used as a proxy for interaction strength, but it is also naturally interpretable as invasion resistance (i.e. whether a rare invader can establish). Framing this more explicitly in terms of invasion fitness would help connect the experimental results to ecological theory and may provide additional insight into the observed patterns.}

\responselabel We agree that Dominance and origin-correlated persistence are not uniquely diagnostic of a single mechanism, and we have revised the manuscript to define community-level selection (CLS) operationally rather than mechanistically. In the revised text, CLS means that persistence during coalescence is correlated with parental-community origin, so that one parental community behaves as a coherent ecological unit; it does not by itself identify whether that coherence arises through direct interspecies competition, shared traits, environmental filtering, or environmental modification by dominant taxa.

This distinction is now made at first use of the term in the Results and expanded in the Discussion, with the framing explicitly linked to prior community-coalescence work \citep{Mansour2018,Rillig2015}. The revised Discussion states that shared environmental tolerances, correlated traits, and pH-mediated filtering by dominant taxa can be mechanisms that generate origin-correlated persistence, rather than mutually exclusive alternatives to it. This is particularly important for Nutr$+$, where dominant-species competition predicts winner identity more strongly in Nutr$+$ than in Base medium ($R^2 = 0.49$ versus $R^2 = 0.11$; Fig.~5C). In the strict acidic--alkaline subset used in Extended Data Fig.~7, where one parental community had pH $<6.5$ and the other had pH $>7.5$, the acidic parental community dominated 29/32 Nutr$+$ events (90.6\%), whereas the Base signal was weaker (23/41 events, 56.1\%). We therefore now describe pH modification as one plausible contributor to winner identity in Nutr$+$, not as a complete explanation for Nutr$+$ Dominance frequency and not as evidence that community-level selection excludes environmental filtering.

At the same time, we do not think a purely passive filtering or hitchhiking explanation is sufficient for the full data set. Several controls narrow that interpretation. Dominance is accompanied by positive within-community pairwise selection correlation in Base and Nutr$+$ ($p < 0.001$ for within-community versus cross-community correlations; Extended Data Fig.~6), whereas Nutr$-$ lacks this correlation. The gLV model, which has identical growth rates and no explicit pH or environmental state variable, reproduces Dominance and positive within-community selection correlation, showing that competitive interaction structure alone is sufficient to generate the outcome-level pattern. In addition, the dominant-species-removal analysis reduces but does not eliminate the Nutr$+$ association between dominant-species competition and PDI (Spearman $\rho = 0.41$--$0.42$, $p \leq 1.3 \times 10^{-3}$; Supplementary Fig.~16), and pre-assembled communities show higher Dominance than directly assembled communities in Base and Nutr$+$ (Extended Data Fig.~8; Supplementary Fig.~23). These results support community-level selection as an outcome-level pattern with multiple mechanistic routes, while acknowledging that the present experiments do not uniquely identify one biochemical mechanism for every event.

We also clarified the ecological meaning of failed invasion. In the revised manuscript, the pairwise invasion assay is described as a direct measure of invasion resistance: an invader introduced at 5\% initial frequency is scored as a failed invasion if it remains below 1\% after seven daily dilution cycles. Failed invasion frequency increases monotonically with nutrient supply (Nutr$-$: $2 \pm 1\%$; Base: $33 \pm 4\%$; Nutr$+$: $48 \pm 4\%$; Fig.~4B), so we now interpret this metric as the empirical proxy for the net intensity of invasion resistance across the nutrient gradient. We further added a Supplementary Note connecting this assay and the pairwise selection-correlation metric to invasion-fitness theory in the gLV framework: invasion fitness is defined as the rare-invader per-capita growth rate in a resident community, and in simulations excess same-parent invasion concordance increases with interaction strength across $\mu = 0.05$--$1.20$ (Pearson $r = 0.870$, $p = 3.2 \times 10^{-8}$; Supplementary Fig.~36). Because our empirical assay begins at 5\% rather than the infinitesimal rare-invader limit, we describe it as an approximation to invasion success and invasion resistance rather than a direct measurement of formal invasion fitness.

\mschange{We revised the Results to define CLS as mechanism-agnostic origin-correlated persistence, revised the Discussion to discuss shared traits, environmental filtering, pH-mediated environmental modification, and hitchhiking as alternative or contributing mechanisms, and added a Supplementary Note connecting failed invasion and pairwise selection correlation to invasion resistance and invasion-fitness theory. The relevant changes appear in Results, ``Community-level selection is prevalent in microbial coalescence'' and ``Theoretical model with random interactions reproduces community-level selection''; Discussion; and Supplementary Note 4 with Supplementary Fig.~36. The experimental definition of failed invasion appears in Methods, ``Pairwise invasion assays.''}
